\section{Методология}

\subsection{HAR-RV модель}

В качестве эконометрического бенчмарка используется модель HAR-RV~\cite{Corsi2009}.
Для повышения стационарности и нормальности остатков оцениваются
логарифмические спецификации:

\textbf{HAR-log}:
\begin{equation}\label{eq:har-log}
    \ln \mathrm{RV}_{t+1} = \beta_0
    + \beta_d \ln \mathrm{RV}_t^{(d)}
    + \beta_w \ln \mathrm{RV}_t^{(w)}
    + \beta_m \ln \mathrm{RV}_t^{(m)}
    + \varepsilon_{t+1},
\end{equation}

\textbf{HAR-log-S} (с sentiment):
\begin{equation}\label{eq:har-s}
    \ln \mathrm{RV}_{t+1} = \beta_0
    + \beta_d \ln \mathrm{RV}_t^{(d)}
    + \beta_w \ln \mathrm{RV}_t^{(w)}
    + \beta_m \ln \mathrm{RV}_t^{(m)}
    + \beta_s S_{t-1}
    + \beta_{sw} \bar{S}_{t-1}^{(w)}
    + \varepsilon_{t+1},
\end{equation}
где $S_{t-1}$~--- лагированный (на один день) sentiment-индекс,
$\bar{S}_{t-1}^{(w)}$~--- его пятидневное скользящее среднее.
Лаг применяется во избежание data leakage: при прогнозе $\mathrm{RV}_{t+1}$
используются только сигналы, известные к концу торгового дня $t-1$.
По умолчанию в качестве $S_{t-1}$ используется взвешенный
sentiment\_combined~(\ref{eq:sent_comb}); в ablation-анализе
(раздел~\ref{sec:har_ablation}) рассмотрены также отдельные источники.

Оценка проводится методом МНК с поправкой HAC (Newey--West, $\ell = 5$ лагов)
для корректного учёта гетероскедастичности и автокорреляции остатков.

\subsection{Схема оценки out-of-sample}

Для обеспечения реалистичной оценки качества прогнозов применяется
процедура расширяющегося окна (expanding window):

\begin{enumerate}
    \item Первые 60\% наблюдений образуют начальную обучающую выборку
          (не менее 120 наблюдений).
    \item На каждом шаге $t$ модель переобучается на всей доступной истории
          $\{1, \ldots, t-1\}$ и даёт прогноз на шаг $t$.
    \item Прогнозы накапливаются и оцениваются совместно.
\end{enumerate}

\subsection{Метрики качества}

\textbf{MAE} (Mean Absolute Error):
\begin{equation}
    \mathrm{MAE} = \frac{1}{T}\sum_{t=1}^{T} |\hat{y}_t - y_t|.
\end{equation}

\textbf{RMSE} (Root Mean Squared Error):
\begin{equation}
    \mathrm{RMSE} = \sqrt{\frac{1}{T}\sum_{t=1}^{T}(\hat{y}_t - y_t)^2}.
\end{equation}

\textbf{QLIKE}~\cite{Patton2011}~--- робастная метрика для волатильности:
\begin{equation}
    \mathrm{QLIKE} = \frac{1}{T}\sum_{t=1}^{T}
    \left(\frac{\mathrm{RV}_t}{\hat{\sigma}_t^2}
    - \ln\frac{\mathrm{RV}_t}{\hat{\sigma}_t^2} - 1\right).
\end{equation}

\textbf{$R^2_{\mathrm{OOS}}$}~\cite{GoyalWelch2008}~---
out-of-sample $R^2$ относительно бенчмарка исторического среднего:
\begin{equation}
    R^2_{\mathrm{OOS}} = 1 - \frac{\mathrm{MSE}_{\text{model}}}{\mathrm{MSE}_{\text{bench}}},
\end{equation}
где $\mathrm{MSE}_{\text{bench}} = T^{-1}\sum_t(y_t - \bar{y})^2$.
Положительный $R^2_{\mathrm{OOS}}$ означает, что модель превосходит
наивный прогноз историческим средним.

\textbf{Тест Диболда--Мариано}~\cite{Diebold1995} применяется для проверки
статистической значимости различий между прогнозами двух моделей.
Для пары моделей $M_1, M_2$ строится ряд разностей функций потерь
$d_t = L(e_t^{(1)}) - L(e_t^{(2)})$ и тестируется
$H_0: \mathbb{E}[d_t] = 0$ против двусторонней альтернативы.
Статистика
\begin{equation}
    \mathrm{DM} = \frac{\bar{d}}{\sqrt{\widehat{\mathrm{Var}}(\bar{d})}}
    \;\xrightarrow{d}\; \mathcal{N}(0, 1),
\end{equation}
где $\widehat{\mathrm{Var}}(\bar{d})$~--- оценка долгосрочной дисперсии
по Newey--West с ядром Бартлетта. В работе тестируются три типа потерь:
SE ($d_t = (y_t-\hat{y}^{(1)}_t)^2 - (y_t-\hat{y}^{(2)}_t)^2$),
AE (абсолютное отклонение) и QLIKE (Patton).

\subsection{XGBoost}

Модель градиентного бустинга XGBoost~\cite{Chen2016}
обучается на расширенном наборе признаков, включающем:
HAR-лаги ($\mathrm{RV}^{(d)}, \mathrm{RV}^{(w)}, \mathrm{RV}^{(m)}$),
дополнительные лаги $\mathrm{RV}_{t-k}$ ($k=2,\ldots,5$),
и sentiment/attention/macro-признаки из \texttt{daily\_sentiment.csv}.
Гиперпараметры: глубина деревьев 4, скорость обучения 0{,}05,
300 деревьев, $L_1$-регуляризация 0{,}1.
Разбиение обучения/теста~--- 60/40 по времени без перемешивания.

\subsection{kNN на пространстве информационных состояний}

По методологии Halousková и Lyócsa~\cite{Halouskova2025}
реализован kNN-прогнозировщик на основе поиска исторически похожих
информационных состояний.
Для каждого тестового дня $t$ строится вектор состояния:
\[
  \mathbf{z}_t = [\underbrace{\mathrm{rv}^{(d)}_t,\, \mathrm{rv}^{(w)}_t,\, \mathrm{rv}^{(m)}_t}_{\text{HAR-блок}},\;
                  \underbrace{s_1,\ldots,s_6}_{\text{sentiment-блок}}],
\]
где $s_1,\ldots,s_6$~--- нормированные sentiment/attention-переменные
(RSS, GDELT, Telegram, Google Trends, объём торгов, combined).
Сходство между днём $t$ и историческим днём $s < t$ вычисляется
как сумма \textit{обратно-расстоянных} (inverse-distance) ядер
по HAR- и sentiment-блокам отдельно:
\[
  \mathrm{sim}(t, s) =
    \frac{1}{1 + \|\mathbf{h}_t - \mathbf{h}_s\|_2} +
    \frac{w_s}{1 + \|\mathbf{s}_{t-1} - \mathbf{s}_{s-1}\|_2},
    \quad w_s = 1{,}0.
\]
Sentiment-вектор $\mathbf{s}$ берётся с лагом~1 (как и в HAR-log-S),
что обеспечивает консистентность всех моделей по графику доступной
информации.
Прогноз~--- взвешенное среднее реализованной волатильности $k$ ближайших соседей:
\[
  \widehat{\ln \mathrm{RV}}_{t+1} = \sum_{i \in \mathcal{N}_k(t)} w_i \cdot \ln \mathrm{RV}_{i+1},
  \quad w_i = \frac{\mathrm{sim}(t,i)}{\sum_{j \in \mathcal{N}_k(t)} \mathrm{sim}(t,j)}.
\]
Оптимальное число соседей $k^*$ подбирается на обучающей выборке.
После устранения leakage и унификации QLIKE-метрики оптимумы
сместились: $k^* = 75$ для GLDRUB\_TOM и $k^* = 30$ для SLVRUB\_TOM.

\subsection{Вычислительная сложность и обоснование оптимальности}
\label{sec:complexity}

Так как проект программный, итоговую модель важно выбирать не только
по точности, но и по вычислительной стоимости в режиме расширяющегося окна.
Пусть $T$~--- число шагов теста, $N$~--- размер обучающей истории на текущем
шаге, $p$~--- число признаков, $m$~--- размерность вектора состояния,
$K$ и $d$~--- число и глубина деревьев бустинга. Сложность моделей приведена
в таблице~\ref{tab:complexity}.

\begin{table}[H]
\centering
\small
\caption{Вычислительная сложность моделей в режиме расширяющегося окна}
\label{tab:complexity}
\begin{tabular}{lll}
\hline
Модель & Переобучение на шаг & Прогноз на шаг \\
\hline
HAR-log / HAR-log-S & $O(N p^2 + p^3)$ & $O(p)$ \\
XGBoost & $O(K\, d\, N \log N)$ & $O(K\, d)$ \\
kNN (информ.\ состояния) & --- (без обучения) & $O(N m)$ \\
\hline
\end{tabular}
\end{table}

HAR-модели сводятся к МНК с малым числом признаков ($p \le 6$), поэтому
переобучение на каждом шаге занимает доли секунды, и суммарная стоимость
$O\!\left(T N p^2\right)$~--- наименьшая среди всех моделей. XGBoost требует
заметно больше вычислений из-за построения деревьев и подбора гиперпараметров.
kNN не обучается, но каждый прогноз требует сравнения со всей историей, что
даёт стоимость $O\!\left(T N m\right)$~--- это самая дорогая модель на этапе
прогноза.

Если сравнить это с результатами раздела~\ref{sec:results}, видно, что
HAR-log-S~--- лучший выбор: она даёт лучшее качество для золота, считается
быстрее всех и сохраняет понятные коэффициенты. Более дорогие XGBoost и kNN
при этом не выигрывают по качеству ни на одном инструменте.
